% 第 7 章：计数原理、概率测度空间与统计线性逼近
\chapter{计数原理、概率测度空间与统计线性逼近}

\section{计数原理的集合本质：分类加法与分步乘法}

\begin{problem}\level{易}
判断下列计数过程分别属于分类加法原理还是分步乘法原理，并计算结果。
(1) 从 4 本不同的语文书与 3 本不同的数学书中任取 1 本，有多少种取法？
(2) 从 4 本不同的语文书与 3 本不同的数学书中各取 1 本，有多少种取法？
(3) 从 4 名男生与 3 名女生中选出 3 人，要求至少有 1 名女生，有多少种选法？
(4) 从 4 名男生与 3 名女生中选出 3 人，分别担任班长、学习委员、生活委员，要求至少有 1 名女生，有多少种选法？
\end{problem}
\begin{solution}
\paragraph{结构剖析}
分类加法原理对应“互斥子集的并集基数”$|A\cup B|=|A|+|B|$，分步乘法原理对应“笛卡尔积的基数”$|A\times B|=|A|\cdot|B|$。判断的关键是看事件是“要么……要么……”（分类，用加法）还是“先……再……”（分步，用乘法）；是否区分顺序则决定用排列 $A_n^m$ 还是组合 $\mathrm{C}_n^m$。
\paragraph{标准解题步骤}
(1) 取语文书或数学书，两类互斥，属分类加法：$4+3=7$ 种。
(2) 先取语文书再取数学书，两步缺一不可，属分步乘法：$4\times3=12$ 种。
(3) “至少有 1 名女生”按女生人数分为 1 女 2 男、2 女 1 男、3 女三类，属分类加法，每类内部用分步乘法：
$\mathrm{C}_3^1\mathrm{C}_4^2+\mathrm{C}_3^2\mathrm{C}_4^1+\mathrm{C}_3^3\mathrm{C}_4^0=3\cdot6+3\cdot4+1=31$ 种；
亦可用间接法 $\mathrm{C}_7^3-\mathrm{C}_4^3=35-4=31$ 验证。
(4) 三个职务有顺序，用排列；先求总排法 $A_7^3=7\cdot6\cdot5=210$，再减去“全为男生”的排法 $A_4^3=4\cdot3\cdot2=24$，得 $210-24=186$ 种。
验算：(1) $7$；(2) $12$；(3) $31$；(4) $186$。
\paragraph{学生典型易错点诊断}
最易犯的错误一是把 (1) 误当分步乘法算成 $4\times3=12$；二是把 (2) 误当分类加法算成 $4+3=7$；三是 (3) 与 (4) 混淆“选人不排序”与“选人再排序”，漏乘职务的排列数；四是 (3) 分类时重复或遗漏（如只算“恰好 1 女”的情形）。规避方法：先问“是否分步”“是否有序”，再决定用加法还是乘法、用排列还是组合。
\end{solution}

\begin{problem}\level{中}
(1) 6 名男生排成一排，4 名女生插入男生之间的空档，要求女生互不相邻，有多少种排法？
(2) 7 名男生排成一排，3 名女生插入空档，要求 3 名女生必须相邻，有多少种排法？
(3) 把 12 个相同的小球放入 3 个不同的盒子，要求每个盒子至少有 2 个小球，有多少种放法？
\end{problem}
\begin{solution}
\paragraph{结构剖析}
插空法处理“不相邻”：先排无限制元素，再在形成的空档中有序地放入被限制元素；捆绑法处理“必须相邻”：把相邻元素视为一个整体参与排列，再乘内部排列；隔板法处理“相同元素分盒且每盒至少若干”：先减掉下界，再用星棒模型计数非负整数解。三者本质都是把“受限排列”翻译为“空档/整体/星棒”的计数。
\paragraph{标准解题步骤}
(1) 先排 6 名男生，有 $6!=720$ 种；6 名男生形成 $7$ 个空档，从中选 4 个并按女生顺序排列，有 $A_7^4=7\cdot6\cdot5\cdot4=840$ 种。由分步乘法得 $720\times840=604800$ 种。
(2) 把 3 名女生捆绑为一个整体，与 7 名男生共 $8$ 个单元全排列，再乘女生内部排列：$8!\cdot3!=40320\times6=241920$ 种。
(3) 设三盒分别放 $x_1,x_2,x_3$ 个，则 $x_i\ge2$ 且 $x_1+x_2+x_3=12$。令 $y_i=x_i-2\ge0$，化为 $y_1+y_2+y_3=6$，非负整数解个数为 $\mathrm{C}_{6+3-1}^{3-1}=\mathrm{C}_8^2=28$ 种。
验算：(1) $604800$；(2) $241920$；(3) $28$。
\paragraph{学生典型易错点诊断}
最易犯的错误一是插空法漏乘空档内女生的排列 $A_7^4$（只选位置不排序）；二是把 6 名男生形成的空档数错记为 6（应为 $6+1=7$）；三是捆绑法忘记乘内部排列 $3!$；四是隔板法忘记“每盒至少 2 个”须先减去 $2\times3=6$，误用 $\mathrm{C}_{11}^2=55$。规避方法：插空数空档为 $n+1$，捆绑必乘内部序，隔板先配方非负。
\end{solution}

\section{概率测度空间：古典概型、几何概型与条件概率}

\begin{problem}\level{易}
袋中有 5 个红球、4 个白球，从中任取 3 个（不放回，且不计顺序）。
(1) 求 3 个全是红球的概率；
(2) 求恰有 1 个白球的概率；
(3) 求至少有 1 个红球的概率。
\end{problem}
\begin{solution}
\paragraph{结构剖析}
古典概型要求样本空间有限且每个基本事件等可能，概率 $P(A)=\dfrac{\text{有利事件数}}{\text{基本事件总数}}$。袋中不放回取球且不计顺序时，样本点用组合计数；关键是把“恰有”“至少有”准确翻译成组合和，并优先用对立事件处理“至少”。
\paragraph{标准解题步骤}
基本事件总数为 $\mathrm{C}_9^3=84$。
(1) 3 个全红：$\mathrm{C}_5^3=10$，故 $P=\dfrac{10}{84}=\dfrac{5}{42}$。
(2) 恰有 1 白即 1 白 2 红：$\mathrm{C}_4^1\mathrm{C}_5^2=4\times10=40$，故 $P=\dfrac{40}{84}=\dfrac{10}{21}$。
(3) “至少有 1 红”的对立事件是“3 个全白”：$\mathrm{C}_4^3=4$，故 $P=1-\dfrac{4}{84}=\dfrac{80}{84}=\dfrac{20}{21}$。
验算：全红 $10$、恰 1 白 $40$、恰 2 白 $\mathrm{C}_4^2\mathrm{C}_5^1=30$、全白 $4$，合计 $84$，概率之和为 $1$。
\paragraph{学生典型易错点诊断}
最易犯的错误一是把“不放回取 3 个”误当有顺序用排列，导致分子分母同乘 $3!$ 后计数混乱；二是“恰有 1 白”误算为 $\mathrm{C}_4^1\mathrm{C}_5^1\mathrm{C}_3^1$；三是“至少有 1 红”直接分类而漏掉情形，应优先用对立事件；四是结果不约分。规避方法：先定“是否有序”，再用组合数，最后用对立事件简化“至少”。
\end{solution}

\begin{problem}\level{中}
(1) 在区间 $[0,6]$ 上随机取两个实数 $x,y$，求 $|x-y|\le2$ 的概率；
(2) 在边长为 $2$ 的正方形内随机投一点，求该点落在以正方形中心为圆心、半径为 $1$ 的圆内的概率；
(3) 甲、乙两人约定在 8 时至 9 时之间随机到达某地，先到者等待 15 分钟后离开，求两人能会面的概率。
\end{problem}
\begin{solution}
\paragraph{结构剖析}
几何概型把概率写成测度之比 $P(A)=\dfrac{\mu(A)}{\mu(\Omega)}$：一维用长度、二维用面积、三维用体积。把随机试验落入坐标系，样本空间是一个区域，有利事件是子区域，面积（长度）之比即概率。
\paragraph{标准解题步骤}
(1) 样本空间 $\Omega=\{(x,y)\mid 0\le x\le6,\ 0\le y\le6\}$，面积 $36$。$|x-y|\le2$ 是夹在两条平行线 $y=x\pm2$ 之间的带状区域，其补集为两个直角边 $4$ 的等腰直角三角形，面积 $2\cdot\dfrac{4^2}{2}=16$，故有利面积 $36-16=20$，$P=\dfrac{20}{36}=\dfrac59$。
(2) $\Omega$ 为边长 $2$ 的正方形，面积 $4$；内切圆半径 $1$，面积 $\pi$，故 $P=\dfrac{\pi}{4}$。
(3) 以甲、乙到达时刻 $x,y\in[0,60]$ 为坐标，$\Omega$ 是边长 $60$ 的正方形，面积 $3600$。会面条件为 $|x-y|\le15$，其补集为两个直角边 $45$ 的等腰直角三角形，面积 $2\cdot\dfrac{45^2}{2}=2025$，有利面积 $3600-2025=1575$，故 $P=\dfrac{1575}{3600}=\dfrac{7}{16}$。
验算：(1) $\dfrac59\approx0.556$；(2) $\dfrac{\pi}{4}\approx0.785$；(3) $\dfrac{7}{16}=0.4375$。
\paragraph{学生典型易错点诊断}
最易犯的错误一是把几何概型误当古典概型去“数点”；二是 (1) 中把带状区域误算成矩形或忘记减去两个三角形；三是 (3) 中把等待时间 $15$ 误当 $30$（两人各等），或把到达区间长度误取为 $60$ 以外的值；四是把面积比写成边长比。规避方法：先画坐标系，标出样本区域与有利区域，再算面积。
\end{solution}

\begin{problem}\level{中}
袋中有 4 个红球、2 个白球，从中不放回地依次取出 2 个。
(1) 已知第一次取到白球，求第二次取到红球的概率；
(2) 求第二次取到红球的概率；
(3) 已知第二次取到红球，求第一次取到白球的概率。
\end{problem}
\begin{solution}
\paragraph{结构剖析}
条件概率 $P(B\mid A)=\dfrac{P(AB)}{P(A)}$ 的本质是把样本空间从 $\Omega$ 缩小到 $A$，并在缩小的空间上重新归一化。不放回抽取时，“已知第一次结果”直接改变剩余球的构成；而“第二次命中”的无条件概率由对称性（或全概率）与第一次无关。
\paragraph{标准解题步骤}
(1) 已知第一次取到白球后，袋中剩 4 红 1 白共 5 个，故 $P=\dfrac45$。
(2) 由对称性，第二次取到红球的概率等于第一次取到红球的概率，即 $\dfrac{4}{6}=\dfrac23$；用全概率验证：$\dfrac{2}{6}\cdot\dfrac45+\dfrac{4}{6}\cdot\dfrac35=\dfrac{8}{30}+\dfrac{12}{30}=\dfrac{20}{30}=\dfrac23$。
(3) $P(\text{第一次白}\mid\text{第二次红})=\dfrac{P(\text{第一次白且第二次红})}{P(\text{第二次红})}=\dfrac{\frac26\cdot\frac45}{\frac23}=\dfrac{\frac{8}{30}}{\frac23}=\dfrac{8}{30}\cdot\dfrac32=\dfrac25$。
验算：(1) $\dfrac45=0.8$；(2) $\dfrac23\approx0.667$；(3) $\dfrac25=0.4$。
\paragraph{学生典型易错点诊断}
最易犯的错误一是把条件概率 $P(B\mid A)$ 与联合概率 $P(AB)$ 混淆；二是 (1) 忘记第一次取出后袋中球数变为 5；三是 (3) 误以为“第二次红球”与第一次无关而直接写 $\dfrac{2}{6}$；四是不验证 (2) 的对称性结论。规避方法：先写出缩小后的样本空间，再写 $P(B\mid A)=\dfrac{P(AB)}{P(A)}$ 并代值。
\end{solution}

\section{全概率公式、贝叶斯公式与离散型分布}

\begin{problem}\level{难}
甲箱中有 3 个红球、3 个白球，乙箱中有 4 个红球、2 个白球。先等可能地任选一个箱子，再从该箱中任取 2 个球。
(1) 求取出的 2 个球恰为 1 红 1 白的概率；
(2) 若取出的 2 个球恰为 1 红 1 白，求它们来自甲箱的概率；
(3) 若取出的 2 个球全是红球，求它们来自甲箱的概率。
\end{problem}
\begin{solution}
\paragraph{结构剖析}
全概率公式把样本空间按“选到哪个箱子”作完备互斥划分 $B_1,B_2$：$P(A)=\sum_i P(B_i)P(A\mid B_i)$；贝叶斯公式是其逆向，用后验概率 $P(B_i\mid A)=\dfrac{P(B_i)P(A\mid B_i)}{P(A)}$ 更新对来源的判断。核心是识别“由因溯果”（全概率）与“由果溯因”（贝叶斯）。
\paragraph{标准解题步骤}
记 $B_1,B_2$ 分别为选到甲、乙箱，$P(B_1)=P(B_2)=\dfrac12$。两箱均有 6 球，取 2 球总数 $\mathrm{C}_6^2=15$。
(1) $P(A\mid B_1)=\dfrac{\mathrm{C}_3^1\mathrm{C}_3^1}{15}=\dfrac9{15}$，$P(A\mid B_2)=\dfrac{\mathrm{C}_4^1\mathrm{C}_2^1}{15}=\dfrac8{15}$，故
$P(A)=\dfrac12\cdot\dfrac9{15}+\dfrac12\cdot\dfrac8{15}=\dfrac{17}{30}$。
(2) $P(B_1\mid A)=\dfrac{\frac12\cdot\frac9{15}}{\frac{17}{30}}=\dfrac{9/30}{17/30}=\dfrac9{17}$。
(3) 记 $C$ 为“2 球全红”：$P(C\mid B_1)=\dfrac{\mathrm{C}_3^2}{15}=\dfrac3{15}$，$P(C\mid B_2)=\dfrac{\mathrm{C}_4^2}{15}=\dfrac6{15}$，$P(C)=\dfrac12\cdot\dfrac3{15}+\dfrac12\cdot\dfrac6{15}=\dfrac{9}{30}=\dfrac3{10}$，故 $P(B_1\mid C)=\dfrac{\frac12\cdot\frac3{15}}{\frac3{10}}=\dfrac{3/30}{3/10}=\dfrac13$。
验算：(1) $\dfrac{17}{30}\approx0.567$；(2) $\dfrac9{17}\approx0.529$（乙箱为 $\dfrac8{17}$）；(3) $\dfrac13\approx0.333$。
\paragraph{学生典型易错点诊断}
最易犯的错误一是把贝叶斯公式的分母写成 $P(A\mid B_1)$ 而非全概率 $P(A)$；二是混淆 $P(B_1\mid A)$ 与 $P(A\mid B_1)$（“由果溯因”与“由因溯果”方向颠倒）；三是两箱红球数不同却误用同一组合数；四是误以为先验概率恒为 $\dfrac12$（若按球数加权则不同）。规避方法：先写划分与先验，再写每个 $P(A\mid B_i)$，最后套贝叶斯公式。
\end{solution}

\begin{problem}\level{中}
某射手每次射击命中目标的概率为 $0.8$，且各次射击相互独立。该射手独立射击 $5$ 次，设命中次数为 $X$。
(1) 求恰好命中 $3$ 次的概率；
(2) 求 $X$ 的数学期望与方差；
(3) 求至少命中 $1$ 次的概率。
\end{problem}
\begin{solution}
\paragraph{结构剖析}
$n$ 重独立重复试验中命中次数 $X\sim B(n,p)$，分布列为 $P(X=k)=\mathrm{C}_n^kp^k(1-p)^{n-k}$，期望 $E(X)=np$，方差 $D(X)=np(1-p)$。期望是“平均命中次数”，方差衡量波动大小；二项分布可看作超几何分布在“有放回/独立”假设下的对应模型。
\paragraph{标准解题步骤}
$X\sim B(5,0.8)$。
(1) $P(X=3)=\mathrm{C}_5^3(0.8)^3(0.2)^2=10\times0.512\times0.04=0.2048$。
(2) $E(X)=np=5\times0.8=4$，$D(X)=np(1-p)=5\times0.8\times0.2=0.8$。
(3) $P(X\ge1)=1-P(X=0)=1-(0.2)^5=1-0.00032=0.99968$。
验算：(1) $0.2048$；(2) $E(X)=4$，$D(X)=0.8$；(3) $0.99968$。
\paragraph{学生典型易错点诊断}
最易犯的错误一是把 $p$ 与 $1-p$ 的位置写反；二是误记 $E(X)=np(1-p)$、$D(X)=np$（期望与方差公式对调）；三是 $\mathrm{C}_5^3$ 误算或漏乘指数项；四是 (3) 直接逐项相加而算错，应优先用对立事件 $1-P(X=0)$。规避方法：先写分布列通项，再代入 $k$，最后用 $E=np$、$D=np(1-p)$ 与对立事件。
\end{solution}

\begin{problem}\level{中}
一批产品共 10 件，其中 3 件是次品、7 件是正品，从中任取 4 件，设次品数为 $X$。
(1) 求恰有 1 件次品的概率；
(2) 求 $X$ 的数学期望与方差；
(3) 求至少有 1 件次品的概率。
\end{problem}
\begin{solution}
\paragraph{结构剖析}
超几何分布描述“从含 $M$ 个次品的 $N$ 个产品中不放回取 $n$ 个，次品数 $X$”的分布：$P(X=k)=\dfrac{\mathrm{C}_M^k\mathrm{C}_{N-M}^{n-k}}{\mathrm{C}_N^n}$，期望 $E(X)=n\dfrac MN$，方差 $D(X)=n\dfrac MN\left(1-\dfrac MN\right)\dfrac{N-n}{N-1}$。与二项分布的区别在于“不放回”，方差多出有限总体修正因子 $\dfrac{N-n}{N-1}$。
\paragraph{标准解题步骤}
$N=10$，$M=3$，$n=4$，$\mathrm{C}_{10}^4=210$。
(1) $P(X=1)=\dfrac{\mathrm{C}_3^1\mathrm{C}_7^3}{\mathrm{C}_{10}^4}=\dfrac{3\times35}{210}=\dfrac{105}{210}=\dfrac12$。
(2) $E(X)=4\times\dfrac{3}{10}=\dfrac65$；$D(X)=4\times\dfrac{3}{10}\times\dfrac{7}{10}\times\dfrac{10-4}{10-1}=\dfrac{14}{25}$。
(3) $P(X\ge1)=1-P(X=0)=1-\dfrac{\mathrm{C}_7^4}{\mathrm{C}_{10}^4}=1-\dfrac{35}{210}=\dfrac56$。
验算：(1) $\dfrac12$；(2) $\dfrac65=1.2$，$\dfrac{14}{25}=0.56$；(3) $\dfrac56\approx0.833$。
\paragraph{学生典型易错点诊断}
最易犯的错误一是把超几何分布误当二项分布，用 $p^k(1-p)^{n-k}$ 计算；二是期望公式写成 $nM$ 或 $\dfrac MN$（漏乘 $n$）；三是方差漏掉有限总体修正因子 $\dfrac{N-n}{N-1}$；四是组合数的上标、下标写反。规避方法：先确认“不放回”，再写超几何公式，最后回代数值检验。
\end{solution}

\section{统计线性逼近：回归直线与独立性检验}

\begin{problem}\level{中}
某班 5 名学生的数学成绩 $x$ 与物理成绩 $y$ 如下表：
\begin{center}
\begin{tabular}{c|ccccc}
$x$ & 1 & 2 & 3 & 4 & 5\\
\hline
$y$ & 2 & 3 & 5 & 4 & 6
\end{tabular}
\end{center}
(1) 求样本中心 $(\bar x,\bar y)$；
(2) 求 $y$ 关于 $x$ 的线性回归方程 $\hat y=\hat b x+\hat a$；
(3) 验证回归直线经过样本中心，并计算残差平方和；
(4) 计算直线 $\hat y=x+1$ 的残差平方和，与 (3) 比较，说明最小二乘法的意义。
\end{problem}
\begin{solution}
\paragraph{结构剖析}
线性回归方程 $\hat y=\hat b x+\hat a$ 是使残差平方和 $Q=\sum(y_i-\hat b x_i-\hat a)^2$ 最小的直线（最小二乘）。对 $Q$ 分别关于 $\hat a,\hat b$ 求偏导并令其为零，可得 $\hat a=\bar y-\hat b\bar x$，故回归直线必过样本中心 $(\bar x,\bar y)$；而 $\hat b=\dfrac{\sum(x_i-\bar x)(y_i-\bar y)}{\sum(x_i-\bar x)^2}$。
\paragraph{标准解题步骤}
(1) $\bar x=\dfrac{1+2+3+4+5}{5}=3$，$\bar y=\dfrac{2+3+5+4+6}{5}=4$，样本中心为 $(3,4)$。
(2) $\sum(x_i-\bar x)(y_i-\bar y)=(-2)(-2)+(-1)(-1)+0\cdot1+1\cdot0+2\cdot2=9$，
$\sum(x_i-\bar x)^2=4+1+0+1+4=10$，故 $\hat b=\dfrac9{10}=0.9$，$\hat a=4-0.9\times3=1.3$，回归方程为 $\hat y=0.9x+1.3$。
(3) 代入 $x=3$：$\hat y=0.9\times3+1.3=4=\bar y$，故直线过样本中心。残差依次为 $-0.2,-0.1,1.0,-0.9,0.2$，残差平方和 $Q=0.04+0.01+1+0.81+0.04=1.9$。
(4) 直线 $\hat y=x+1$ 的残差为 $0,0,1,-1,0$，残差平方和 $Q'=0+0+1+1+0=2>1.9$，可见回归直线确实使残差平方和最小。
验算：(2) $\hat b=0.9$，$\hat a=1.3$；(3) $Q=1.9$；(4) $Q'=2$。
\paragraph{学生典型易错点诊断}
最易犯的错误一是把 $\hat b$ 的分子、分母写反；二是求 $\hat a$ 时误用 $\hat a=\bar y+\hat b\bar x$（符号错误）；三是计算 $\sum x_iy_i$ 时漏项；四是误以为回归直线是“过最多点”的直线，而它只保证残差平方和最小；五是把残差与残差平方和混淆。规避方法：先算样本中心，再用偏差公式求 $\hat b$，最后验证过中心并比较 $Q$。
\end{solution}

\begin{problem}\level{中}
某校为研究“是否选考物理”与“是否参加竞赛培训”的关系，随机调查了 160 名学生，得到 $2\times2$ 列联表如下：
\begin{center}
\begin{tabular}{c|cc|c}
 & 选考物理 & 不选考物理 & 合计\\
\hline
参加竞赛培训 & 50 & 30 & 80\\
未参加竞赛培训 & 30 & 50 & 80\\
\hline
合计 & 80 & 80 & 160
\end{tabular}
\end{center}
(1) 计算 $K^2$ 的观测值；
(2) 已知临界值满足 $P(K^2\ge7.879)=0.005$，$P(K^2\ge10.828)=0.001$，判断是否有 $99.5\%$ 的把握认为两者有关；
(3) 简述独立性检验所体现的“小概率反证法”逻辑。
\end{problem}
\begin{solution}
\paragraph{结构剖析}
独立性检验把“两个分类变量是否有关”转化为对 $2\times2$ 列联表计算
$K^2=\dfrac{n(ad-bc)^2}{(a+b)(c+d)(a+c)(b+d)}$，再与临界值比较。其逻辑是统计学反证法：先假设“无关”，若在该假设下观测到 $K^2$ 很大（对应小概率事件），则有理由拒绝原假设。
\paragraph{标准解题步骤}
由列联表读出 $a=50$，$b=30$，$c=30$，$d=50$，$n=160$。
(1) $ad-bc=50\times50-30\times30=1600$，故
$$K^2=\frac{160\times1600^2}{80\times80\times80\times80}=\frac{160\times2\,560\,000}{40\,960\,000}=10.$$
(2) 因为 $7.879<10<10.828$，所以有 $99.5\%$ 的把握认为“是否选考物理”与“是否参加竞赛培训”有关，但没有 $99.9\%$ 的把握。
(3) 先假设两者无关，此时 $K^2$ 应较小；观测值 $10$ 远大于临界值 $7.879$，说明在“无关”假设下出现该数据是小概率事件，故拒绝原假设，认为两者有关。
验算：$K^2=10$；临界值 $7.879$（$\alpha=0.005$）、$10.828$（$\alpha=0.001$）。
\paragraph{学生典型易错点诊断}
最易犯的错误一是把公式分母 $(a+b)(c+d)(a+c)(b+d)$ 写错或漏项；二是把 $n$ 只乘分子的一部分；三是把“有 $99.5\%$ 把握”误表述为“有关的概率是 $99.5\%$”；四是忘记先假设无关（反证法的前提）。规避方法：先画列联表定位 $a,b,c,d,n$，再套公式，最后按临界值表下结论。
\end{solution}
